\documentclass[UTF8,zihao=-4]{ctexart}
\usepackage[a4paper,margin=2.2cm]{geometry}
\usepackage{xcolor}
\usepackage{hyperref}
\usepackage{enumitem}
\usepackage{longtable}
\usepackage{array}
\hypersetup{colorlinks=true,linkcolor=teal,urlcolor=teal}
\setlist{itemsep=0.25em,topsep=0.35em}
\title{\bfseries 点子发芽日报 2026-05-24}
\author{点子发芽}
\date{2026-05-24}
\begin{document}
\maketitle
\section*{今日总结}
今天的输入共同指向一个轻量雷达：一边是 Codex goal 这类工具能力，一边是 AI 升学/项目机会和线下商业观察。它们不必被硬凑成同一主题，但都值得用“关键词、补充信息、权重”的方式先收进来，再判断哪些能变成行动线索。
\section*{今日输入}
\begin{longtable}{|p{0.22\linewidth}|p{0.58\linewidth}|p{0.10\linewidth}|}
\hline
\textbf{关键词} & \textbf{补充信息} & \textbf{权重} \\
\hline
codex新增/goal功能 & 公众号推送，后来又上B站学习了一下相关知识 & 3 \\
\hline
深圳科创学院 & 学长分享未来升学路径，AI相关 & 3 \\
\hline
上海创智学院 & 学长分享未来升学路径，AI相关 & 3 \\
\hline
私域 & 听0.2\%创业交流群里一个开雪茄吧的人说的 & 3 \\
\hline
\end{longtable}
\section*{今日新知}
\subsection*{codex新增/goal功能}
\paragraph{简介} Codex 的 goal 功能可以理解为一种面向持续任务的工作模式：用户先给出一个明确目标，Codex 围绕这个目标反复读取文件、运行命令、修改方案并验证结果，直到目标完成或遇到真实阻塞。它和普通问答的区别在于，普通问答更像一次响应，goal 更像把一段工作流托管给助手；它强调目标、约束、执行顺序、验证和收尾汇报。适合处理跨多个文件、多轮调整、需要检查结果的任务，例如重构提示词、修复日报质量、补齐脚本和文档。
\paragraph{最近有什么相关新闻} 本次可用来源显示，Codex 相关文档已经把围绕目标持续推进作为明确用例，并同时提供提示词和命令协作说明 [S1] [S2] [S3]。这些资料更适合作为功能核实依据，不宜扩写成产品发布新闻。
\paragraph{与我相关} 补充信息显示你先从公众号被动看到，再主动去 B 站补课，说明这个词不是一闪而过的热词，而是已经触发了工具使用层面的学习。它和当前项目的关系很直接：报告质量修复、跨天自动化、补跑验证这类任务都适合用 goal 模式持续推进。
\paragraph{最小下一步} 记录一次使用 goal 模式完成的真实任务：目标、输入材料、验证命令、最终产物各写一行。
\subsection*{深圳科创学院}
\paragraph{简介} 深圳科创学院是围绕硬科技创新、创业训练和工程人才培养建立的机构线索。它的关键词不只是“学院”，更重要的是“科创”：通常会涉及真实问题、技术原型、创业营、导师网络、产业资源和团队孵化。理解它时应先把它放在高校课堂之外的实践平台来看，而不是简单等同于传统升学院校。对这类机构，真正需要确认的是参与门槛、训练形式、项目质量、导师背景、与高校或企业的连接，以及它能否沉淀为项目经历或研究方向。
\paragraph{最近有什么相关新闻} 当前搜索到的资料主要来自机构官网和高校通知，能说明它长期围绕科创训练、创业营和学生实践机会组织活动 [S6] [S7] [S8]。这些来源适合做入口核实，但还需要继续追问具体项目质量。
\paragraph{与我相关} 补充信息来自学长分享的升学路径，而且明确和 AI 相关，这说明你关心的不是机构名本身，而是它能不能提供 AI+工程、AI+硬件、项目制训练或创业实践的机会。下一步应把学长口中的经验拆成可验证问题，而不是只收藏名称。
\paragraph{最小下一步} 问学长一个具体问题：深圳科创学院最有价值的是项目经历、导师网络、创业训练，还是升学申请叙事。
\subsection*{上海创智学院}
\paragraph{简介} 上海创智学院是一个面向人工智能人才培养、科研转化和产学研协同的新型机构线索。理解它时，重点不在“学院”这个名称，而在它是否提供 AI 相关课程、项目、导师、实习、成果转化或跨校资源。与传统院系相比，这类机构往往更强调产业牵引、真实场景、交叉学科和高水平平台连接。判断它的价值时，需要区分宣传口径和可参与入口：哪些活动面向本科生，哪些需要研究基础，哪些只是资讯，哪些能实际进入项目。
\paragraph{最近有什么相关新闻} 公开资料显示，它围绕 AI 人才培养生态和成果发布持续有机构信息更新，官网和高校/科委渠道可作为优先核实来源 [S11] [S12] [S13]。这些信息能证明它是可追踪对象，但不能替代对参与路径的具体确认。
\paragraph{与我相关} 补充信息同样来自学长关于未来升学路径的分享，所以它应和深圳科创学院放在同一组比较，而不是孤立成一个知识节点。你真正需要的是判断它是否能连接 AI+X、实习、科研项目、导师资源或申请叙事。
\paragraph{最小下一步} 做一张两行小表，比较上海创智学院和深圳科创学院：适合对象、可参与项目、能否接触导师或实习。
\subsection*{私域}
\paragraph{简介} 私域指的是品牌、商家或个人能够相对自主、反复触达并持续经营的一组用户关系。它不是某个单一工具，也不等于微信群、企微好友或会员系统；这些只是承载关系的渠道。私域的核心是把一次性流量沉淀成可识别、可沟通、可复购、可转介绍的关系资产。判断一个私域案例是否成立，要看用户为什么愿意留下、商家如何持续提供价值、触达是否被用户接受，以及这种关系最终是否提升复购、客单价、活动参与或信任度。
\paragraph{最近有什么相关新闻} 这次检索到的多为机构营销页和行业文章，适合作为普通线索，不适合作为强事实依据 [S16] [S17] [S18]。可保留的判断是：高客单、低频、强信任的线下店更适合观察私域如何从关系经营中产生复购。
\paragraph{与我相关} 补充信息来自创业交流群里的雪茄吧案例，这比抽象定义更有价值。雪茄吧通常涉及高客单、线下社交、熟人信任和活动组织，正好可以观察私域如何从一次消费变成会员关系、圈层连接和反复到店。
\paragraph{最小下一步} 把雪茄吧案例补成三行：顾客为什么留下、老板如何触达老客、复购或活动靠什么发生。
\section*{与旧知识的链接}
\begin{itemize}[leftmargin=2em]
\item Codex goal 和本地节点“新知孵化工作流”直接相关：它可能降低跨天维护报告和自动化任务时反复交代背景的成本。
\item 深圳科创学院、上海创智学院可以合并成一个观察主题：AI+X 升学/实习机会雷达，先比较机会类型，不急着判断高低。
\item 私域可以暂时作为商业观察素材，而不是独立升成高权重节点；它和雪茄吧案例绑定后，才有个人辨识度。
\end{itemize}
\section*{今日发芽点子}
\begin{itemize}[leftmargin=2em]
\item AI+X 升学/实习机会雷达：用一张小表持续比较不同学院、训练营和项目入口。
\item 高客单线下店的私域机制观察：从雪茄吧开始积累真实商业案例。
\end{itemize}
\section*{参考搜索内容}
\begin{itemize}[leftmargin=2em]
\item \href{https://developers.openai.com/codex/use-cases/follow-goals}{[S1] 官方/机构资料：Follow a goal | Codex use cases - developers.openai.com}
\item \href{https://developers.openai.com/codex/prompting}{[S2] 官方/机构资料：Prompting - Codex | OpenAI Developers}
\item \href{https://developers.openai.com/codex/app/commands}{[S3] 官方/机构资料：Commands - Codex app | OpenAI Developers}
\item \href{https://innoxsz.com/explore-innox}{[S6] 官方/机构资料：走进学院}
\item \href{http://due.hitsz.edu.cn/info/1486/6195.htm}{[S7] 官方/机构资料：关于开展2025深圳科创学院冬季科创营报名的通知-哈尔滨工业大学(深圳)教务部}
\item \href{https://www.innoxsz.com/}{[S8] 官方/机构资料：深圳科创学院}
\item \href{https://www.ecnu.edu.cn/info/1426/70270.htm}{[S11] 官方/机构资料：成立一年来,上海创智学院打破传统教育框架,重构顶尖ai人才培养生态 全年猎才破边界 课堂跟着产业跑}
\item \href{https://www.sast.gov.cn/content.html?id=1949032550202621954&cid=36}{[S12] 官方/机构资料：破界而来！上海创智学院首发八大AI成果}
\item \href{https://www.sii.edu.cn/}{[S13] 官方/机构资料：上海创智学院}
\item \href{https://www.weiling.cn/info/tMMWt3aZ.html}{[S16] 项目主页/机构主页：提升私域高客单高复购与流量管理的三大策略-卫瓴科技官网}
\item \href{https://www.vzkoo.com/read/20250702636339fd321e457645743369.html}{[S17] 普通线索：2023年实体零售私域运营分析：门店私域流量构建已成破局关键}
\item \href{https://www.sohu.com/a/917190938_120405735}{[S18] 普通线索：揭秘高客单行业的私域运营策略\_用户\_品牌\_社群 - 搜狐}
\end{itemize}
\end{document}
